折射率 \(n>1\) 且忽略色散的透明介质以匀速 \(v\) 沿光传播轴向运动，\(|v|\ll c\)。介质自身静止系内光的相速度大小是 \(c/n\)，光在实验室系传播路径长度 \(L\) 固定。两束光分别顺介质运动和逆介质运动通过等长路径后合束干涉，光真空波长为 \(\lambda\)。按狭义相对论速度合成，只保留 \(v/c\) 一阶。
(1) 从 Lorentz 速度合成推导实验室中顺、逆方向的光速一阶式。
(2) 求两束光通过实验室固定路径的时间差。
(3) 推导干涉相位差和条纹移动数。
(4) 检验 (n	o1) 与 (v	o0) 极限，指出忽略色散时适用范围。
